% Generated draft: requires mathematical and source review.
\documentclass[11pt]{article}
\usepackage[margin=25mm]{geometry}
\usepackage{fontspec}
\setmainfont{Latin Modern Roman}
\usepackage{amsmath,amssymb,mathtools,unicode-math}
\setmathfont{Latin Modern Math}
\usepackage{enumitem,needspace,xurl,hyperref}
\setlength{\parindent}{0pt}
\setlength{\parskip}{5pt}
\setlength{\emergencystretch}{3em}
\newcommand{\N}{\mathbb{N}}
\newcommand{\Z}{\mathbb{Z}}
\newcommand{\Q}{\mathbb{Q}}
\newcommand{\R}{\mathbb{R}}
\newcommand{\C}{\mathbb{C}}
\newcommand{\F}{\mathbb{F}}
\newcommand{\A}{\mathbb{A}}
\newcommand{\PP}{\mathbb{P}}
\newcommand{\E}{\mathbb{E}}
\newcommand{\eps}{\varepsilon}
\newcommand{\dd}{\,\mathrm{d}}
\begin{document}
\section*{UnsolvedMath ID: 5600009}
\textbf{Source problem code:} catalogue entry
\subsection{Definitions and mathematical statement}
Let $n\ge1$ and let $H_1,\ldots,H_{n+2}$ be projective hyperplanes in $\PP^n(\C)$ in general position. Thus homogeneous defining linear forms for any $n+1$ of them are linearly independent; in particular no $n+1$ of these hyperplanes meet. Put
\[
 B=\bigcup_{j=1}^{n+2}H_j.
\]
Here subtraction of the source's set of hyperplanes means subtraction of their union.

The supplied higher-dimensional Picard assertion is that no nondegenerate holomorphic map
\[
 f:\C^n\longrightarrow\PP^n(\C)\setminus B
\]
exists. Under the usual maximal-rank meaning for maps between equidimensional complex manifolds, this can be written
\[
 \forall f:\C^n\to\PP^n(\C)\text{ holomorphic},\qquad
 f(\C^n)\cap B=\varnothing
 \ \Longrightarrow\
 \operatorname{rank}_{\C}(\mathrm df_z)<n
 \quad\text{for every }z\in\C^n.
\]
The record does not define nondegenerate. Linear nondegeneracy, meaning the image is not contained in any projective hyperplane, is a different condition. The intended condition must be confirmed rather than treating these two meanings as equivalent. The displayed rank formulation is explicitly conditional on the maximal-rank reading.
\subsection{Short English statement}
Can a nondegenerate holomorphic map from complex $n$-space omit $n+2$ hyperplanes in general position? The supplied assertion says it cannot. The registry labels this solved, and the research summary attributes an impossibility result to Carlson--Griffiths-type defect relations or a linear-degeneracy theorem, without exact references or hypotheses.

The source is Chern's 1970 presentation, with the title attributing the conjecture to Wu. Provenance records only its historical open status and says current status needs review. The later solved assessment is retained with the unresolved meaning of nondegeneracy and its missing theorem citation.
\subsection{Sources}
\begin{itemize}
\item Local UnsolvedMath registry, permanent ID 5600009, source code catalogue entry, ``Wu's Higher-Dimensional Picard Conjecture,'' listed as proposed by Shiing-Shen Chern, 1970. The title's Wu attribution is retained without inferring an original bibliographic source. Dataset curation and metadata: Ulam AI and UnsolvedMath Contributors, CC BY 4.0.
\item Shiing-Shen Chern, \emph{Open Problems in Differential Geometry} (1970), Problem 9, PDF page 10, public ICM-proceedings scan, \url{https://people.math.harvard.edu/~hirolee/pdfs/2014-fall-230a-icm1970-chern-differential-geometry.pdf}; normalized pdf-text extraction, accessed 2026-07-29.
\item The research summary cites Carlson--Griffiths-type defect relations or a linear-degeneracy theorem without a title, date, precise theorem statement, or publication locator.
\item Source-specific provenance: the historical item is stated open, current status NEEDS\_REVIEW, and statement-reuse terms for the public ICM scan NEEDS\_REVIEW. These qualifications are separate from the dataset licence.
\end{itemize}
\end{document}
